	\section{Méthode de gestion de projet : l'approche Kanban}
\label{sec:kanban}
%=============================================================

La conduite d'un projet informatique de la nature de celui présenté
dans ce travail alliant traitement de données statistiques,
développement d'une application web requiert une méthode de gestion de projet adaptée
à la complexité, à l'incertitude et à l'évolutivité inhérentes à ce
type de réalisation. Parmi les différentes approches disponibles,
la méthode \textbf{Kanban} a été retenue comme cadre organisationnel
du présent projet, en raison de sa flexibilité, de sa capacité à
visualiser en temps réel l'état d'avancement des tâches et de son
adéquation avec les contraintes d'un projet conduit par un nombre
limité de collaborateurs. La présente section expose les fondements
théoriques et pratiques de cette méthode, avant d'en présenter
l'application concrète dans le cadre du développement de la
plateforme \textit{Business Partnership}.

%-------------------------------------------------------------
\subsection{Origines et fondements conceptuels de la méthode Kanban}
\label{subsec:kanban_origines}
%-------------------------------------------------------------

La méthode Kanban trouve ses origines dans le système de production
industrielle développé par l'ingénieur Taiichi Ohno au sein des
usines Toyota dans le Japon de l'après-guerre, au milieu des années
1950. Inspiré du fonctionnement des supermarchés américains, dans
lesquels les rayons sont réapprovisionnés en fonction de la demande
réelle des consommateurs et non selon des prévisions préétablies,
Ohno conçoit un système de production dit \textit{en flux tiré}
(\textit{pull system}), dans lequel chaque étape du processus de
fabrication ne produit que ce que l'étape suivante est en mesure
d'absorber \cite{ohno_1988}. Ce principe fondamental de régulation
par la demande constitue le c\oe ur philosophique de la méthode
Kanban, dont le nom signifie littéralement \og{}étiquette\fg{} ou
\og{}signal visuel\fg{} en japonais, en référence aux cartes physiques
utilisées pour matérialiser le flux de travail sur les lignes
d'assemblage.

Transposée au domaine du développement logiciel et de la gestion de
projets informatiques au début des années 2000, la méthode Kanban
a été formalisée et popularisée par David J. Anderson, dont l'ouvrage
de référence \textit{Kanban: Successful Evolutionary Change for Your
	Technology Business} (2010) constitue la base théorique sur laquelle
repose la quasi-totalité des implémentations modernes de cette
méthode \cite{anderson2015essential}. Anderson identifie cinq propriétés
fondamentales qui caractérisent toute implémentation Kanban~:

\begin{enumerate}
	\item \textbf{La visualisation du flux de travail}~: l'ensemble des
	tâches constituant le projet est rendu visible sur un tableau
	structuré en colonnes représentant les différents états du processus
	de réalisation (à faire, en cours, terminé, etc.), permettant à
	chaque membre de l'équipe d'appréhender instantanément l'état
	d'avancement global du projet~;
	
	\item \textbf{La limitation du travail en cours (WIP)}~: chaque
	colonne du tableau Kanban est soumise à une limite maximale de
	tâches simultanément en cours (\textit{Work In Progress limit},
	ou WIP limit), afin d'éviter la surcharge des ressources, de
	favoriser l'achèvement des tâches engagées avant d'en démarrer
	de nouvelles, et d'identifier rapidement les goulots d'étranglement
	du processus~;
	
	\item \textbf{La gestion et la mesure du flux}~: le flux de travail
	est mesuré et géré de manière continue, à l'aide d'indicateurs
	tels que le \textit{lead time} (délai entre la création d'une tâche
	et sa livraison) et le \textit{cycle time} (délai de traitement
	effectif), afin d'identifier les inefficiences et d'optimiser
	la cadence de livraison~;
	
	\item \textbf{La formalisation des règles du processus}~: les
	critères d'entrée et de sortie de chaque étape du processus sont
	explicitement définis et partagés par l'ensemble des parties
	prenantes, garantissant ainsi la cohérence et la prévisibilité
	des livraisons~;
	
	\item \textbf{L'amélioration continue collaborative}~: la méthode
	Kanban inscrit l'équipe dans une démarche permanente d'amélioration
	du processus (\textit{kaizen}), fondée sur l'analyse régulière
	des métriques de flux et la mise en \oe uvre collective de mesures
	correctives.
\end{enumerate}

Kanban s'inscrit ainsi dans la famille des méthodes agiles de gestion
de projet, dont le Manifeste Agile de 2001 a posé les principes
fondateurs~: priorité aux individus et aux interactions sur les
processus, aux logiciels fonctionnels sur la documentation exhaustive,
à la collaboration avec le client sur la négociation contractuelle,
et à l'adaptation au changement sur le suivi d'un plan rigide
\cite{beck_2001}. Il se distingue cependant d'autres méthodes agiles
telles que Scrum en ce qu'il ne prescrit ni itérations de durée fixe
(\textit{sprints}), ni rôles prédéfinis, ni cérémonies obligatoires,
offrant ainsi une flexibilité organisationnelle maximale, particulièrement
adaptée aux projets conduits par des équipes de petite taille ou en
contexte académique.

%-------------------------------------------------------------
\subsection{Structure et composantes du tableau Kanban}
\label{subsec:kanban_tableau}
%-------------------------------------------------------------

Le tableau Kanban constitue l'artefact central de la méthode~: il
matérialise visuellement l'ensemble du flux de travail et en rend
l'état instantanément lisible par l'ensemble des parties prenantes
du projet. Dans sa forme la plus courante, le tableau Kanban est
organisé en colonnes verticales représentant les différentes étapes
du cycle de vie d'une tâche, et en lignes horizontales pouvant
correspondre à différentes catégories de travaux ou à différents
membres de l'équipe. Chaque tâche est représentée par une carte
(\textit{kanban card}), physique ou numérique, qui se déplace de
gauche à droite sur le tableau au fur et à mesure de son avancement.

Dans le cadre du développement de la plateforme \textit{Business
	Partnership}, le tableau Kanban a été structuré autour de
\textbf{cinq colonnes} correspondant aux étapes successives du
processus de développement~:

\begin{enumerate}
	\item \textbf{Backlog} (\textit{À planifier})~: cette colonne
	regroupe l'ensemble des tâches identifiées mais non encore
	planifiées pour réalisation. Elle constitue le réservoir de travail
	du projet et est alimentée au fur et à mesure de l'identification
	de nouveaux besoins fonctionnels ou techniques~;
	
	\item \textbf{À faire} (\textit{To Do})~: les tâches sélectionnées
	pour être traitées dans la période courante sont déplacées du
	Backlog vers cette colonne. Leur sélection est effectuée en
	fonction de leur priorité et de la capacité disponible de l'équipe~;
	
	\item \textbf{En cours} (\textit{In Progress})~: cette colonne
	regroupe les tâches en cours de réalisation. Une limite WIP est
	appliquée à cette colonne afin de garantir la focalisation de
	l'effort sur un nombre restreint de tâches simultanées~;
	
	\item \textbf{En révision} (\textit{In Review})~: les tâches
	achevées sont soumises à une phase de relecture, de test et de
	validation avant d'être considérées comme définitivement terminées.
	Cette étape garantit la qualité des livrables produits~;
	
	\item \textbf{Terminé} (\textit{Done})~: les tâches ayant satisfait
	l'ensemble des critères de validation définis sont archivées dans
	cette colonne, constituant ainsi un historique des réalisations
	accomplies.
\end{enumerate}

\usetikzlibrary{shadows} % <--- Indispensable pour l'effet "drop shadow"

\begin{figure}[H]
	\centering
	\begin{tikzpicture}[
		col/.style={rectangle, draw=primary, fill=primary!8,
			minimum width=2.2cm, minimum height=7cm,
			text width=2cm, align=center, rounded corners=3pt},
		card/.style={rectangle, draw=secondary, fill=white,
			minimum width=1.8cm, minimum height=0.8cm,
			text width=1.7cm, align=center, font=\tiny,
			rounded corners=2pt, drop shadow},
		header/.style={rectangle, fill=primary, text=white,
			minimum width=2.2cm, minimum height=0.6cm,
			text width=2cm, align=center, font=\small\bfseries,
			rounded corners=3pt}
		]
		% Colonnes
		\foreach \i/\name in {0/{Backlog}, 2.4/{À faire}, 4.8/{En cours}, 7.2/{En révision}, 9.6/{Terminé}}{
			\node[col] at (\i, 0) {};
			\node[header] at (\i, 3.2) {\name};
		}
		% Cartes exemples
		\node[card] at (0, 2.2) {Conception MCD};
		\node[card] at (0, 1.2) {Tests unitaires};
		\node[card] at (0, 0.2) {Documentation};
		\node[card] at (2.4, 2.2) {ACM / CAH};
		\node[card] at (2.4, 1.2) {Module auth};
		\node[card, fill=secondary!15] at (4.8, 2.2) {Algo matching};
		\node[card, fill=secondary!15] at (4.8, 1.2) {Interface profil};
		\node[card, fill=accent!10] at (7.2, 2.2) {API REST};
		\node[card, fill=green!20] at (9.6, 2.2) {Base de données};
		\node[card, fill=green!20] at (9.6, 1.2) {Page d'accueil};
		% Limite WIP
		\node[font=\tiny\color{accent}] at (4.8, -3) {WIP max : 2};
	\end{tikzpicture}
	\caption{Représentation schématique du tableau Kanban utilisé pour
		le développement de la plateforme \textit{Business Partnership}}
	\label{fig:kanban_board}
\end{figure}
%-------------------------------------------------------------
\subsection{Avantages de Kanban pour un projet de Data Sciences}
\label{subsec:kanban_avantages}
%-------------------------------------------------------------

Le choix de la méthode Kanban comme cadre de gestion du présent
projet se justifie par plusieurs avantages spécifiques, particulièrement
pertinents dans le contexte d'un projet de Data Sciences combinant
analyse statistique et développement logiciel.

\textbf{La flexibilité face à l'incertitude analytique} constitue
le premier avantage déterminant. Un projet de Data Sciences est par
nature itératif et exploratoire~: les résultats d'une analyse
statistique peuvent remettre en question des choix de modélisation
antérieurs ou faire émerger de nouveaux besoins fonctionnels pour
la plateforme. Contrairement à des méthodes prescriptives telles
que le cycle en V ou le modèle en cascade, Kanban permet d'intégrer
ces évolutions sans rupture organisationnelle, simplement en ajoutant
ou en repriorisant des cartes dans le tableau. Cette adaptabilité
est d'autant plus précieuse que les phases d'analyse (ACM, CAH, AFD)
et de développement (Django, SQLite) s'influencent mutuellement de
manière non linéaire au cours du projet.

\textbf{La visibilité permanente de l'avancement} représente le
second avantage majeur. Le tableau Kanban offre à tout moment une
photographie fidèle de l'état du projet~: les tâches terminées, en
cours, en attente et à planifier sont visibles d'un seul regard,
sans nécessiter de réunions de suivi formelles ou de rapports
d'avancement élaborés. Cette transparence est particulièrement utile
dans le contexte d'un projet académique impliquant des interactions
régulières avec un directeur de mémoire et un encadreur professionnel,
à qui l'état d'avancement peut être communiqué de manière instantanée
et sans ambiguïté.

\textbf{La prévention des goulots d'étranglement} constitue le
troisième avantage opérationnel. La limitation du WIP contraint à
achever les tâches engagées avant d'en démarrer de nouvelles, ce
qui prévient l'accumulation de travaux partiellement réalisés et
favorise une livraison régulière et prévisible de fonctionnalités.
Dans le cadre du développement de la plateforme \textit{Business
	Partnership}, cette discipline a permis d'éviter la dispersion entre
les différentes composantes du projet (analyse des données,
développement du back-end, développement du front-end, rédaction du
mémoire) et de maintenir une progression équilibrée sur l'ensemble
des axes.

\textbf{La légèreté procédurale}, enfin, constitue un avantage
non négligeable pour un projet conduit en contexte académique. Kanban
ne requiert ni réunion quotidienne obligatoire (\textit{daily
	stand-up}), ni définition de rôles spécialisés (\textit{product
	owner}, \textit{scrum master}), ni planification d'itérations à
durée fixe, ce qui le rend directement applicable sans formation
préalable approfondie ni infrastructure organisationnelle complexe.
Sa mise en \oe uvre peut être réalisée à partir d'outils simples et
accessibles, aussi bien physiques (tableau blanc et post-its) que
numériques, comme les plateformes Trello, Notion ou GitHub Projects,
qui ont été mobilisées dans le cadre de ce projet.

%-------------------------------------------------------------
\subsection{Application de Kanban au développement de la plateforme
	\textit{Business Partnership}}
\label{subsec:kanban_application}
%-------------------------------------------------------------

L'application concrète de la méthode Kanban au développement de la
plateforme \textit{Business Partnership} s'est traduite par
l'organisation de l'ensemble des tâches du projet en une liste de
travaux (\textit{backlog}) initiale, établie dès la phase de
conception et enrichie de manière incrémentale tout au long du
développement. Cette liste a été structurée en trois catégories
principales de tâches~: les tâches d'analyse statistique (traitement
des données de l'enquête, réalisation de l'ACM, de la CAH et de
l'AFD), les tâches de développement logiciel (conception de la base
de données, développement des modèles Django, développement des vues
et des templates, intégration de l'algorithme de recommandation), et
les tâches de documentation et de rédaction (rédaction du mémoire,
documentation technique de la plateforme).

La priorisation des tâches a été conduite selon le principe de la
valeur métier~: les fonctionnalités à plus forte valeur pour
l'utilisateur final (inscription des entreprises, saisie des besoins,
affichage des recommandations) ont été traitées en priorité, afin
de disposer le plus tôt possible d'un prototype fonctionnel
susceptible d'être présenté et évalué. Cette approche, cohérente
avec les principes agiles de livraison incrémentale de valeur,
a permis de valider rapidement les choix de conception et d'ajuster
le développement en fonction des retours obtenus lors des revues
intermédiaires avec l'encadrement académique et professionnel.

Le tableau~\ref{tab:kanban_taches} ci-dessous présente une sélection
des principales tâches réalisées au cours du projet, organisées par
catégorie et accompagnées de leur état d'avancement à la date de
rédaction de ce mémoire.

\begin{table}[H]
	\centering
	\caption{Sélection des principales tâches du projet \textit{Business Partnership} dans le tableau Kanban}
	\label{tab:kanban_taches}
	\renewcommand{\arraystretch}{1.3}
	\begin{tabularx}{\textwidth}{l X >{\centering\arraybackslash}p{2.5cm}}
		\toprule
		\rowcolor{primary!15}
		\textbf{Catégorie} & \textbf{Tâche} & \textbf{État} \\
		\midrule
		
		\multirow{4}{*}{\textbf{Analyse}}
		& Prétraitement des données de l'enquête & \textcolor{green!60!black}{Terminé} \\
		& Réalisation de l'ACM (FactoMineR) & \textcolor{green!60!black}{Terminé} \\
		& Réalisation de la CAH (méthode de Ward) & \textcolor{green!60!black}{Terminé} \\
		& Réalisation de l'AFD et calcul du TBC & \textcolor{green!60!black}{Terminé} \\
		
		\midrule
		
		\multirow{5}{*}{\textbf{Développement}}
		& Conception du MCD et du MLD (MERISE) & \textcolor{green!60!black}{Terminé} \\
		& Développement des modèles Django & \textcolor{green!60!black}{Terminé} \\
		& Développement du module d'inscription & \textcolor{green!60!black}{Terminé} \\
		& Intégration de l'algorithme de recommandation & \textcolor{green!60!black}{Terminé} \\
		& Tests et validation de la plateforme & \textcolor{orange}{\textbf{En cours}} \\
		
		\midrule
		
		\multirow{3}{*}{\textbf{Rédaction}}
		& Rédaction des chapitres 1 à 3 & \textcolor{green!60!black}{Terminé} \\
		& Rédaction du chapitre 4 & \textcolor{orange}{\textbf{En cours}} \\
		& Relecture et mise en forme finale & \textcolor{blue}{\textbf{À faire}} \\
		
		\bottomrule
	\end{tabularx}
\end{table}
L'utilisation de la méthode Kanban tout au long de ce projet a
permis de maintenir une organisation rigoureuse et transparente,
de livrer de manière incrémentale les différentes composantes de
la plateforme et de l'analyse, et d'adapter continuellement les
priorités en fonction de l'évolution des contraintes académiques
et techniques. Elle constitue en ce sens un cadre méthodologique
pleinement cohérent avec la nature exploratoire et itérative d'un
projet de Data Sciences appliquées au développement d'une plateforme
gouvernementale.
\section{Structure du mémoire}

Ce document est organisé en quatre chapitres. Le \textbf{Chapitre~1} présente la structure d'accueil et le déroulement du stage. Le \textbf{Chapitre~2} propose une revue de la littérature et un état de l'art sur les plateformes similaires. Le \textbf{Chapitre~3} décrit la méthodologie adoptée pour la conception et le développement de la plateforme. Enfin, le \textbf{Chapitre~4} présente et discute les résultats obtenus.
